% Chapter 10: The Turn
% Scene function: Sunday evening. Gordon reads Middlemarch aloud.
% Tert watches from the hallway. The irony stops being enough.
% The most restrained chapter in the book. Every line carries weight.

\chapter{Middlemarch}

Sunday evening.

I am in the hallway, positioned for the night shift. Gordon has eaten dinner, washed the dishes, watched television until nine-thirty. The routine has not changed in three weeks. What has changed is the quality of his attention: he moves through the house more carefully now, checking rooms before he enters them, pausing at thresholds. This is my work. I know it is my work because I built the conditions that produced it. Or it is his grief. I have no way to tell.

At nine-forty he turns off the television. He does not go to the bedroom. He goes to the bookshelf.

He stands in front of it for a moment. Then he reaches up to the third shelf on the right and takes down the thick paperback that has been lying flat on top of the others since before I arrived. The one with the cracked spine. The one I noted in the survey photograph and again on my first visit and dismissed as not relevant to the work.

The cover says \emph{Middlemarch.} I have seen it on the shelf for three weeks without reading the title.

He carries the book to his armchair and sits down. He opens it to the bookmark, about two-thirds through. He reads for a moment in silence, finding his place.

Then he begins to read aloud.

Gordon has not read aloud on any previous evening I have observed. This is a new pattern. I note it the way I note all new patterns, as data, as a behavioral event to be logged in the margin pad.

But his voice is different when he reads aloud. I have been listening to Gordon for three weeks: the way he talks to himself while cooking, the way he answers the phone, the way he says goodnight to no one before he turns off the bedroom light. This is not those voices. This is the voice of a man who spent thirty years reading to rooms full of teenagers who did not want to be there and who learned to make them want to be there. It has warmth and pacing and a specific attention to the weight of each sentence, the care of a voice being used for someone else's benefit.

I catch a fragment of what he is reading: ``\ldots the element of tragedy which lies in the very fact of frequency\ldots'' I am not listening to the words. I am listening to the voice.

He is not reading to the room. He is reading to someone who is not here, in the voice he would use if she were, and the distinction between those two things is the largest thing I have encountered in this house.

He reads for about forty minutes. I do not time it exactly. I am not, for these forty minutes, timing anything.

He pauses between chapters. During the pauses he does not put the book down. He holds it open, as if the conversation is only resting.

When he finishes he puts the bookmark back, closes the book, and places it flat on top of the others on the third shelf. This is how it looked in the survey photograph. This is how it has looked every time I have been in this room.

He turns off the lamp. He goes to bed. He closes the bedroom door most of the way but not all the way, the way a person does when they have spent years listening for sounds from another room.

I stay in the hallway for a while.
